%!TEX root = ../main.tex

\section{Основные соотношения}

Пусть плотность внутренней энергии среды $e = \hat{e}(\eta, \vD, \phi, \nabla \phi, \rho)$. Плотность свободной энергии $\psi$ получается из $e$ преобразованием Лежандра относительно пар переменных $\theta$--$\eta$ и $\vE$--$\vD$ и, таким образом,
\begin{align}
    \psi &= e - \vE \cdot \vD - \theta \eta,
    \label{eq:free_energy_density} \\
    \psi &= \hat{\psi}(\theta, \vE, \phi, \nabla \phi, \rho).
    \nonumber
\end{align}

Уравнение баланса микросил и микронапряжений имеет вид
\begin{equation}
    \Div \vxi + \pi + \gamma = 0.
    \label{eq:microforces}
\end{equation}

Закон сохранения внутренней энергии принимается в виде
\begin{equation}
    \dt{e} = -\Div \vq - \Div[\vE, \vect{H}] + \Div(\dt{\phi} \vxi) + \gamma \dt{\phi} - \Div(\kappa \rho \vj) + r,
    \label{eq:energy_conservation}
\end{equation}
где $\vq$ -- плотность теплового потока, $[\vE, \vect{H}]$ -- плотность потока энергии электромагнитного поля, $\Div(\dt{\phi} \vxi)$ -- работа микронапряжений, $\gamma \dt{\phi}$ -- работа внешних микросил, $r$ -- работа внешних сил. Вектор $\kappa \rho \vj$, по-видимому, выражает поток энергии, связанный с диффузией заряда; он добавлен в уравнение, чтобы при дальнейшем выводе получить нужное определяющее соотношение для тока $\vj$ и положительный диапазон температуры $\theta$.

В качестве энтропийного неравенства постулируется
\begin{equation}
    \dt{\eta} \geqslant -\Div \left( \cfrac{\vq}{\theta} \right) + \cfrac{r}{\theta}.
    \label{eq:entropy_inequality}
\end{equation}